\documentclass[11pt,a4paper]{article}
%% Shared preamble: fonts, listings, theorem environments.
\usepackage{fontspec}
\usepackage{amsmath,amssymb,amsthm}
\usepackage{listings}
\usepackage{xcolor}
\usepackage[margin=1.1in]{geometry}
\usepackage{hyperref}
\hypersetup{colorlinks=true, linkcolor=blue!50!black, urlcolor=blue!50!black,
            citecolor=blue!50!black}

%% DejaVu covers most of the Lean symbol set; STIX Two Math fills the rest.
\directlua{luaotfload.add_fallback("monofb",
  {"file:STIXTwoMath-Regular.otf:mode=node;",
   "file:DejaVuSans.ttf:mode=node;"})}
\setmonofont{DejaVu Sans Mono}[Scale=MatchLowercase,
  RawFeature={fallback=monofb}]

\lstdefinelanguage{lean}{
  morekeywords={theorem,lemma,def,structure,inductive,mutual,end,namespace,
    instance,abbrev,where,match,with,fun,let,have,show,from,by,intro,exact,
    refine,cases,rcases,obtain,induction,simp,omega,rfl,deriving,open,import,
    variable,attribute,example,at,do,if,then,else,noncomputable},
  sensitive=true,
  morecomment=[l]{--},
  morecomment=[n]{/-}{-/},
  morestring=[b]",
}
\lstset{
  language=lean,
  basicstyle=\ttfamily\footnotesize,
  keywordstyle=\color{blue!55!black}\bfseries,
  commentstyle=\color{green!35!black}\itshape,
  stringstyle=\color{red!50!black},
  columns=flexible,
  breaklines=true,
  breakatwhitespace=true,
  showstringspaces=false,
  frame=leftline,
  framesep=6pt,
  xleftmargin=10pt,
  numbers=none,
  captionpos=b,
  %% one fullwidth capital appears in a docstring and is in neither font
  literate={Ｄ}{{D}}1,
}

\theoremstyle{plain}
\newtheorem{theorem}{Theorem}[section]
\newtheorem{lemma}[theorem]{Lemma}
\newtheorem{proposition}[theorem]{Proposition}
\theoremstyle{definition}
\newtheorem{definition}[theorem]{Definition}
\newtheorem{remark}[theorem]{Remark}
\newtheorem{finding}[theorem]{Finding}


\title{A Forward Refutation Calculus for Intuitionistic Logic,\\
       Mechanised --- and the First Step Toward the Lax Modality}
\author{Matthew Fairtlough \quad (with Claude Opus 5)}
\date{16 August 2026}

\begin{document}
\maketitle

\begin{abstract}
This is a record of a mechanisation, in Lean 4, of Fiorentini and
Ferrari's forward refutation calculus \textbf{FRJ}$(G)$ for intuitionistic
propositional logic, and of the first step of its extension with the lax
modality~$\bigcirc$ of Propositional Lax Logic. The mechanised base is
sorry-free and choice-free: soundness, completeness and the biconditional
between refutability and the existence of a countermodel are all proved,
the completeness proof \emph{computes} a refutation from a countermodel,
and the only appeal to classical logic in the whole development is
isolated in one statement, where it is a property of the statement and not
of the proof. The extension step adds the modality to the syntax and to
the semantics while leaving every result about the modality-free fragment
intact; two things the build discovered along the way are recorded as
findings, one of them a correction to the plan the extension was following.
All Lean code appears in Appendix~\ref{app:code}, and the exposition
quotes from it throughout.
\end{abstract}

\tableofcontents

%% ==================================================================
\section{Provenance, and what is borrowed}
\label{sec:prov}

The calculus is due to Camillo Fiorentini and Mauro Ferrari, \emph{Duality
between Unprovability and Provability in Forward Refutation-search for
IPL}, ACM Transactions on Computational Logic 21(3), Article 22, 2020.
\textbf{The definitions, the rule table, the statements of the numbered
results and the structure of their proofs are theirs}; wording is borrowed
from that article throughout, and where a definition or a proof outline
below reads like the original that is because it is the original,
transcribed. What is ours is the mechanisation, the divergences recorded
as divergences, and everything from \S\ref{sec:w1} onward.

Two texts carry this calculus and they are not the same text. The
mechanisation was transcribed from the \LaTeX{} source of
arXiv:1804.06689, and cross-read against the published article. Inside the
material formalised here the published version

\begin{itemize}
\item adds a lemma, that every member of the stable set of a derivable
  irregular sequent is smaller than the sequent's right formula, whose
  proof uses one of the refutation-search restrictions;
\item states part (ii) of the key soundness lemma with the hypothesis
  ``the p-sequent forces the whole stable set'' along a restricted
  reachability relation, where the arXiv text has ``forces the stable set
  intersected with the proper subformulas of the goal'' along the plain
  one;
\item swaps the numbering of two properties inside the join case, renames
  the search restrictions, and moves the height bounds to a later section.
\end{itemize}

The arXiv form of part (ii) is the \emph{stronger} statement --- weaker
hypothesis, larger range of quantification --- and it needs none of the
search restrictions, so it is the one mechanised. Numbering below is the
published article's. The correspondence, established by reading both:

\begin{center}
\begin{tabular}{lll}
\hline
Result & journal & arXiv preprint\\
\hline
Soundness of the calculus & Theorem 3.1 & Theorem 1\\
Left formulas along the step relation & Lemma 3.5 & Lemma 1\\
The key soundness lemma & Lemma 3.10 & Lemma 2\\
The extracted model is a countermodel & Theorem 3.12 & Theorem 2\\
Rank equals extracted height & Lemma 6.2 & Lemma 13\\
The completeness construction & Lemma 6.3 & Lemma 14\\
Minimal-height refutation & Theorem 6.4 & Theorem 11\\
The closure of the determining part & Lemma 6.7 & Lemma 15\\
Completeness & Theorem 5.13(i) & Theorem 10\\
\hline
\end{tabular}
\end{center}

\paragraph{Notation.} The article writes the two zones of an irregular
sequent with Greek capitals that Lean reserves. Throughout this paper and
the development they are written \texttt{St} (the \emph{stable} set, the
article's own word for it) and \texttt{Th}. Contexts of regular sequents
keep the letter $\Gamma$.

%% ==================================================================
\section{Why a refutation calculus}
\label{sec:why}

Under the standing discipline of this project a claim is \textsc{proved}
only when it is machine-checked, and otherwise \textsc{refuted} --- by a
kernel-checked countermodel --- or \textsc{open}. If a refutation costs
orders of magnitude more than a proof, every campaign accumulates open
cells it cannot afford to settle. The existing decision procedure for
Propositional Lax Logic gives an effective \emph{proof} procedure but an
exhaustive generate-and-test \emph{disproof} procedure. A refutation
calculus reads a countermodel off a derivation directly, which is what
makes disproof as cheap as proof.

FRJ$(G)$ is that calculus for intuitionistic logic, and it has two
properties that make it the right base. It is \emph{forward}: derivations
are built from axioms upward, over a finite set of sequents fixed by the
goal formula~$G$, so termination is structural. And its soundness is a
\emph{model construction}: a derivation is not merely a certificate that a
countermodel exists, it \emph{is} one, once read correctly.

%% ==================================================================
\section{Preliminaries}
\label{sec:prelim}

\subsection{Language}

The language has propositional variables, the connectives $\land$, $\lor$,
$\supset$ and the constant $\bot$; negation is defined. To it we add the
unary modality $\bigcirc$, present in the syntax from the start of
\S\ref{sec:w1} and mentioned by no rule.

\begin{lstlisting}
inductive Form where
  | atom : String → Form
  | bot : Form
  | and : Form → Form → Form
  | or : Form → Form → Form
  | imp : Form → Form → Form
  | circ : Form → Form
  deriving DecidableEq, Repr
\end{lstlisting}

\noindent
The full text is Appendix~\ref{app:basic}. The size of a formula is its
number of symbols; the modality contributes one.

\subsection{Left and right subformulas}

The calculus is goal-directed: every formula occurring in a sequent is a
subformula of $G$, and \emph{which side} of a sequent it may occur on is
fixed by the polarity of its occurrence in $G$. The article defines the
sets of left and right subformulas as the smallest subsets of the
subformulas of $G$ closed under: $G$ is a right subformula; a conjunction
or disjunction passes its polarity to both components; an implication in
left position puts its consequent left and its antecedent right, and
dually. The modality passes its polarity to its argument, like conjunction
and disjunction and unlike implication --- forced, since $\bigcirc$ is
monotone in its argument.

Rather than define the two sets inductively and then compute them, the
development computes them by a mutual recursion, and \emph{proves} the
article's closure clauses of the computed pair. That proof is the fidelity
check on the definition:

\begin{lstlisting}
structure SfClosed (R L : List Form) : Prop where
  rAnd : ∀ {A B : Form}, Form.and A B ∈ R → A ∈ R ∧ B ∈ R
  rOr  : ∀ {A B : Form}, Form.or A B ∈ R → A ∈ R ∧ B ∈ R
  rImp : ∀ {A B : Form}, Form.imp A B ∈ R → A ∈ L ∧ B ∈ R
  rCirc : ∀ {A : Form}, Form.circ A ∈ R → A ∈ R
  lAnd : ∀ {A B : Form}, Form.and A B ∈ L → A ∈ L ∧ B ∈ L
  lOr  : ∀ {A B : Form}, Form.or A B ∈ L → A ∈ L ∧ B ∈ L
  lImp : ∀ {A B : Form}, Form.imp A B ∈ L → A ∈ R ∧ B ∈ L
  lCirc : ∀ {A : Form}, Form.circ A ∈ L → A ∈ L
\end{lstlisting}

\subsection{Models and forcing}

A model is a finite poset with a minimum and a monotone valuation. Two
divergences from the article, both deliberate and both recorded in the
source (Appendix~\ref{app:basic}).

First, finiteness is carried as a \emph{constructive enumeration} ---~a
list of all worlds together with a proof that it is complete, plus
decidable equality, order and valuation~--- rather than as a finiteness
\emph{property}. The reason is not aesthetic: eliminating a finiteness
property into data requires a choice principle, and the development is
choice-free. The payoff is immediate: forcing becomes a
\emph{decision procedure}, not merely a proposition.

Second, the clause for implication is written as an implication rather
than as the article's disjunction ``$\beta \nVdash A$ or $\beta \Vdash
B$''; writing the disjunction would put excluded middle into the
definition itself.

\begin{lstlisting}
def force (K : Kripke) : K.W → Form → Prop
  | _, .bot => False
  | a, .atom p => K.V a p
  | a, .and A B => force K a A ∧ force K a B
  | a, .or A B => force K a A ∨ force K a B
  | a, .imp A B => ∀ b, K.le a b → force K b A → force K b B
  | a, .circ A => ∀ b, K.le a b → ∃ c, K.Rm b c ∧ force K c A
\end{lstlisting}

\noindent
The last clause and the relation \texttt{Rm} it consults are the subject
of \S\ref{sec:w1}; the five above it are the article's.

\subsection{The closure}

The single most important auxiliary notion is the \emph{closure} of a set
of formulas, generated by the grammar
\[
X \;::=\; C \mid X \land X \mid A \lor X \mid X \lor A \mid A \supset X
\qquad (C \in \Gamma,\ A \text{ any formula}).
\]
It is the set of formulas that a world forces \emph{merely because} it
forces $\Gamma$, and it is what allows the calculus to have no left rules
at all: conjunction and disjunction on the left are absorbed into it. Six
properties are proved (Appendix~\ref{app:basic}); the first, that forcing
$\Gamma$ implies forcing its closure, is the one the rules consume, and
the fifth, that the closure adds no new propositional variable, is what
makes the valuation of the extracted model well defined.

Note where the modality is \emph{not}: the grammar has no clause for it.
That absence is deliberate and is discussed in
Finding~\ref{find:zone}.

%% ==================================================================
\section{The calculus}
\label{sec:calc}

\subsection{Sequents}

There are two kinds. A \emph{regular} sequent $\Gamma \Rightarrow C$ is
read ``some world forces $\Gamma$ and refutes $C$''. An \emph{irregular}
sequent $\texttt{St} ; \texttt{Th} \to C$ carries two zones: the stable
set, which must be preserved when the sequent is used, and a second zone
some of whose members may be lost. Contexts are drawn from
\[
\hat G \;=\; \hat G_{\mathrm{at}} \cup \hat G_{\supset},
\]
the propositional variables and the implications among the left
subformulas of $G$ --- and the fact that these two exhaust the possible
context members is exactly what conjunction and disjunction being absorbed
by the closure buys.

\subsection{The rules}

Two axioms, right-introduction rules for the three connectives, and two
\emph{join} rules; there are no left rules. The two joins are the only way
to obtain a regular sequent from irregular ones, and they are where the
countermodel grows: the conclusion of a join corresponds to a \emph{new
world}, lying below the worlds of its premises.

The join rules take $n \geq 1$ irregular premises. Writing the zones of
the $j$-th premise split into their atomic and implicational parts, the
conclusion keeps the union of the stable atomic parts, the intersection of
the second zones' atomic parts, the union of the stable implicational
parts, and the intersection of the second zones' implicational parts
\emph{restricted} to those implications whose antecedent is the right
formula of some premise. Two side conditions govern them: each premise's
stable set is contained in the union of every other premise's two zones,
and every implication in a stable set is \emph{supported}, meaning its
antecedent is the right formula of some premise.

The support condition is the heart of the calculus and deserves a
sentence of semantics. The new world must force the implications kept in
its context. An implication is forced at a world when its antecedent fails
there; support arranges exactly that, by making the antecedent the goal of
a premise, and the goal of a premise is what fails at the new world.

The Lean rule table is Appendix~\ref{app:calculus}. Each constructor
carries every side condition as a field, and the figure's blanket
condition, that the right formula of every conclusion is a right
subformula of $G$, is a field on every constructor.

\subsection{Canonical contexts}
\label{sec:nf}

One design point in the mechanisation has no counterpart in the article
and is worth stating because it is load-bearing.

Contexts denote sets. On the article's finite sets, the split of a zone as
$\texttt{Th} \cup \Lambda$ demanded by the irregular implication rule is a
literal equality. On lists it is not: concatenation is neither commutative
nor idempotent. Nor can one transport along ``same members'', because
\emph{same members does not imply same derivations} --- the irregular
axiom pins its own zone. The resolution is a change of carrier, not of
calculus: represent every computed context canonically, as the members of
$\hat G$ lying in it.

\begin{lstlisting}
def nf (G : Form) (l : List Form) : List Form :=
  (gHat G).filter (fun x => decide (x ∈ l))

theorem nf_ext {G : Form} {l m : List Form}
    (h : ∀ x, x ∈ gHat G → (x ∈ l ↔ x ∈ m)) : nf G l = nf G m
\end{lstlisting}

\noindent
Two contexts with the same members inside $\hat G$ are then \emph{literally
the same list}: the property that concatenation lacks and finite sets had.
This is legitimate exactly because every context of a derivable sequent is
a subset of $\hat G$ --- a lemma, proved in Appendix~\ref{app:step} ---
and because every side condition in the rule table is membership-based.

%% ==================================================================
\section{Soundness}
\label{sec:sound}

Soundness is the statement that a refutation of $G$ yields a countermodel
for $G$, and its proof is a construction.

\subsection{The step relation}

Write $\sigma_1 \mapsto_0 \sigma_2$ when $\sigma_2$ is the conclusion of a
rule and $\sigma_1$ one of its premises, and $\mapsto$ for the transitive
closure. The first lemma controls how the left formulas move along it.

\begin{lemma}[Lemma 3.5]
\label{lem:lhs}
If $\sigma_1 \mapsto_0 \sigma_2$ by a rule other than the
implication-not-in rule then the left formulas of $\sigma_2$ are among
those of $\sigma_1$; in general they lie in the closure of those of
$\sigma_1$; and the same holds along the transitive closure.
\end{lemma}

This is the lemma that makes the extracted model's valuation monotone, via
the closure property that no new variable is created.

\subsection{The extracted model}

Call a regular sequent of a derivation a \emph{p-sequent} when it is an
axiom or the conclusion of a join. The worlds of the extracted model are
the p-sequents; one lies below another when the second reaches the first
by the step relation; and the valuation of a world is the set of
propositional variables among its left formulas.

The mechanisation builds this in pieces rather than all at once, because
the recursion on derivations is what produces it. The key construction
places a fresh root below the disjoint union of a family of models:

\begin{lstlisting}
def addRoot {ι : Type} [DecidableEq ι] (ιe : List ι) (ιc : ∀ i, i ∈ ιe)
    (V₀ : List Form) (M : ι → Kripke)
    (hV : ∀ (i : ι) (p : String), Form.atom p ∈ V₀ → (M i).V (M i).root p) :
    Kripke where ...
\end{lstlisting}

\noindent
and the lemma that makes working in pieces legitimate is that forcing is
unchanged at a component world (Appendix~\ref{app:model}). This is not a
reformulation of the article's model: a join's conclusion is a fresh
p-sequent lying below exactly the p-sequents of its premises, and every
p-sequent has a unique p-sequent immediately below it, so the article's
order \emph{is} computed by this placement.

\subsection{The key lemma}

\begin{lemma}[Lemma 3.10]
\label{lem:key}
For every sequent occurring in a refutation:
(i) if it is regular, the p-sequent immediately above it forces its
context and refutes its goal;
(ii) if it is irregular, then any p-sequent it reaches which forces its
stable set intersected with the proper subformulas of its goal refutes
that goal.
\end{lemma}

The proof is a main induction on the height of the sequent in the
derivation, by cases on the last rule. The join case carries a
\emph{secondary} induction, on the size of a formula, proving two things
simultaneously: that every implication kept in the conclusion is forced at
the new world, and that the right formula of every premise fails there.
The two feed each other --- the first at an implication calls the second
at its antecedent, and the second calls the first at the members of a
stable set that are proper subformulas of the goal --- and it is the
strict decrease of size that makes the mutual call well founded. This is
the only place in the development where the article's proof has real
depth, and the mechanisation follows it step for step; the article's own
proof is in its appendix and was read there.

\begin{theorem}[Theorems 3.12 and 3.1]
The extracted model is a countermodel for $G$; hence if $G$ is refutable
then $G$ is not valid.
\end{theorem}

%% ==================================================================
\section{Completeness}
\label{sec:complete}

The article proves completeness twice: once in the main text, as a
corollary of a duality with a second calculus and the correctness of a
search procedure, and once in its minimality section, directly, by a
construction on an arbitrary countermodel. The second route is
mechanised. It is equally published, it is a fraction of the work, and ---
decisively --- it \emph{computes}.

\subsection{The determining part}

Fix a countermodel and a world $\alpha$. Say that $\alpha$ forces a
formula \emph{determinately} when it forces it and, in the case of an
implication, fails its antecedent. The \emph{determining part} at $\alpha$
is the set of left subformulas of $G$ determinately forced there. The
point of it is:

\begin{lemma}[Lemma 6.7, in the two directions actually used]
\label{lem:det}
Everything in the determining part is forced at $\alpha$; and every left
subformula of $G$ forced at $\alpha$ lies in the closure of the
determining part.
\end{lemma}

\begin{remark}
The article states this as a chain of set equalities, and taken literally
the first of them is false: the closure grammar lets its implication
clause range over \emph{all} antecedents, so the closure of the
determining part contains implications that need not be left subformulas
of $G$. The proof also has a step that silently needs a conjunction to be
a left subformula. The two directions the construction actually uses are
true, and are what is proved. This divergence is recorded in the
development's fidelity table.
\end{remark}

\subsection{The construction}

\begin{lemma}[Lemma 6.3]
\label{lem:minmod}
For every world $\alpha$ and every right subformula $C$ of $G$ not forced
at $\alpha$ there are: an irregular refutation whose stable set lies
between the determining part and its union with the second zone; and a
regular refutation whose context contains the determining part of some
world above $\alpha$.
\end{lemma}

The proof is a triple induction --- on the height of the world, then on
the kind of sequent with irregular before regular, then on the size of the
goal --- and in the mechanisation it is a \emph{function}, returning the
refutation, because extracting a refutation from an existence statement
needs a choice principle. That is the whole point:

\begin{lstlisting}
def completenessData {G : Form} (hcf : ∀ X ∈ sfR G ++ sfL G, X.isCirc = false)
    (K : Kripke) (hK : ¬ K.valid G) : Derivation G :=
  let w := minMod K G hcf K.root 1 G (sfR_self G) hK
  ⟨w.ctx, w.der⟩
\end{lstlisting}

\noindent
an algorithm from countermodel to refutation. The side condition
\texttt{hcf} is the subject of \S\ref{sec:w1}; on a modality-free goal it
holds automatically, and before \S\ref{sec:w1} it was not there at all.

\subsection{Two simplifications}

The article's construction carries rank bounds, which serve its
minimal-height theorem and are not needed for completeness; and it obeys
four refutation-search restrictions, which exist to shrink the search
space. Both are dropped. Dropping restrictions can only make soundness
\emph{stronger} --- more derivations to be sound about --- and costs
completeness nothing, because the construction happens to satisfy them
anyway.

%% ==================================================================
\section{Choice, and where it is unavoidable}
\label{sec:choice}

The development is choice-free, and this was not free. Two sources had to
be removed. The finite-set library of the ambient mathematical library
reports a choice axiom at the level of \emph{definitions}: union, deletion
and image all do, so a term that merely mentions them carries choice
however it is proved. The list operations do not. And one automation
tactic reasons classically. So: lists throughout, shape predicates as
booleans, a constructive enumeration in place of a finiteness property,
and explicit proofs where the tactic was used.

What remains is a genuine divergence, and it is a property of a
\emph{statement}. The passage from ``not every model validates $G$'' to
``some model refutes $G$'' is not constructively valid. The development
therefore states the biconditional twice:

\begin{lstlisting}
theorem frj_iff_countermodel (G : Form) (hcf : ...) :
    Provable G ↔ ∃ K : Kripke, ¬ K.valid G          -- choice-free

theorem frj_iff_not_IPL (G : Form) (hcf : ...) :
    Provable G ↔ ¬ IPL G                            -- the article's
\end{lstlisting}

\noindent
and the second is the only place in the whole development where a choice
axiom appears. The axiom audit (Appendix~\ref{app:audit}) pins this, and
every pin is guarded so that a regression is a build failure rather than a
discovery months later.

\textbf{Consequence for anything modal}: results about the modality belong
on the countermodel side of that line.

%% ==================================================================
\section{The modality: the first step}
\label{sec:w1}

The extension is staged. The first stage adds the modality to the syntax
and to the semantics, adds no rule, and requires that every
modality-free result still holds. The point of doing it this way is that
the eight modules then say, by breaking or not breaking, exactly which
proofs are modality-sensitive.

\subsection{What was added}

The formula constructor, its clauses in size, subformulas and polarity;
a second accessibility relation on models, a preorder contained in the
first; the forcing clause
\[
K,\alpha \Vdash \bigcirc A \quad\text{iff}\quad
\text{for every } \beta \geq \alpha \text{ there is } \gamma
\text{ with } \texttt{Rm}\,\beta\,\gamma \text{ and } K,\gamma \Vdash A,
\]
and the corresponding clause of the decision procedure for forcing, so
that forcing stays decidable --- with no axioms at all. The universal
quantifier over $\beta$ is what makes the modality monotone without any
interaction axiom between the two frames.

The result: \texttt{lake build FRJ} is green, and the axiom audit file is
\emph{untouched} --- every pin still passes.

\subsection{Two findings}

\begin{finding}[The fresh-root construction needs its own modal relation]
\label{find:arrm}
The obvious choice, taking the modal relation of the assembled model to be
its order, makes the component-preservation lemma \textbf{false} for the
modality: at a component world the modal clause would then quantify over
the component's \emph{order}-successors, where the component itself
quantifies over its own modal successors. The assembled model must
therefore carry a separate inductive relation which keeps each component's
modal relation and lets the fresh root see everything modally.
\end{finding}

\begin{lstlisting}
inductive ARRm {ι : Type} (M : ι → Kripke) :
    Option ((i : ι) × (M i).W) → Option ((i : ι) × (M i).W) → Prop
  | root (x) : ARRm M none x
  | comp {i a b} : (M i).Rm a b → ARRm M (some ⟨i, a⟩) (some ⟨i, b⟩)
\end{lstlisting}

\begin{finding}[The third zone cannot precede the rules]
\label{find:zone}
The modality needs a third context zone. Semantically: the closure absorbs
conjunction and disjunction on the left, which is why two zones are
exhaustive for intuitionistic logic, and it absorbs the modality no
better than it absorbs implication, since $\bigcirc A$ can be forced where
$A$ is not, exactly as $A \supset B$ can be forced where $B$ is not.
Empirically, from an earlier screen over 156 cells: with the determining
part taken to be variables, falsehood and implications there are 32
certified failures; adding the modal formulas gives none.

\emph{Nevertheless the third zone cannot be added at this stage}, and the
build is what showed it. With a third zone present, the side condition of
the implication-not-in rule admits a modal formula into a zone; the
irregular implication rule can then shift it into a stable set; and the
join rules, which keep only the atomic and implicational parts, silently
\emph{drop} it --- breaking the condition that the join case of
Lemma~\ref{lem:key} has to verify. The third zone, the modal clause of the
closure and the modal clause of the determining part are therefore
\emph{one change together with the rules}, not a change before them.
\end{finding}

\subsection{The asymmetry}

Soundness came through the extension \textbf{unchanged}: it needs no side
condition, because it never asks what determines a world. Completeness
gained one --- that every subformula of the goal is modality-free ---
threaded from Lemma~\ref{lem:det} through the construction to
\texttt{completenessData} and the two biconditionals. That is the honest
statement of what the modality-free theory proves: it is a theory of
modality-free goals, and the determining part is where the modality first
bites.

%% ==================================================================
\section{The modal rules: screens, and the design}
\label{sec:rules}

The rules are new mathematics: the calculus has a published source, its
modal extension does not. They are therefore screened before anything is
proved about them, and the screens are part of the library
(Appendix~\ref{app:modallean}): five semantic obligations proved with no
axioms, and every cell settled by the decision procedure for forcing.

\subsection{The difficulty, in one sentence}

For implication the obligation at a new world is \emph{negative} --- the
antecedent fails there, so the implication holds vacuously --- which is
all the support condition has to arrange, by naming the antecedent as some
premise's goal. For the modality the obligation is \emph{positive}: a
modal witness must exist, and nothing in the calculus's data supplies one.

\begin{lstlisting}
theorem circ_intro {w : K.W} {A : Form}
    (wit : ∃ u, K.Rm w u ∧ K.force u A)
    (above : ∀ v, K.le w v → v ≠ w → K.force v (.circ A)) :
    K.force w (.circ A)

theorem not_force_circ {w : K.W} {A : Form}
    (h : ∀ u, K.Rm w u → ¬ K.force u A) : ¬ K.force w (.circ A)
\end{lstlisting}

\subsection{Three screens}

\paragraph{The witness must be selective.} In the two-world model with the
variable true only at the upper world and the modal relation equal to the
order, the root forces $\bigcirc p$ and refutes $p$. So a rule must be able
to build a world whose modal successor forces something the world itself
does not.

\paragraph{The witness is per-world.} In the three-world model with a root
below two incomparable worlds carrying $p$ and $q$, the root forces
$\bigcirc(p \lor q)$ while both $\bigcirc p$ and $\bigcirc q$ fail. So the
modal data cannot be one promise fixed for the whole derivation: it
belongs to the world a join creates. The same model refutes the
distribution of the modality over disjunction.

\paragraph{The witness cannot be an arbitrary context.} This screen is a
refutation of a design rather than a model. The tempting rule declares
``there is a world forcing $\Delta$'' for a context of its choosing. That
is exactly the design refuted in the abandoned first attempt at this
extension, whose world predicate constrained a zone only by membership in
the universe, with no closure condition; three kernel-checked cells killed
it. So a declared context carries a \emph{realisability} obligation, and
realisability is not membership. The obligation is stated in the library
and nothing proves it, which is the point.

\subsection{The zones arrive with the rules}

Finding~\ref{find:zone} makes the following one \emph{atomic} change, each
part of which is unsound without the others: the context universe gains
its third zone; the closure gains a modal clause; the determining part
gains modal formulas; the join rules gain a modal zone and its support
condition; and the modal rules are added. Because canonical contexts are
computed by filtering against the universe, the first part
re-canonicalises \emph{every} context in the development --- which is the
cost of the step, and the reason it must be done once rather than in
pieces.

\subsection{The rules}

\paragraph{Regular modal introduction, gated by barrenness.}
\[
\frac{\Gamma \Rightarrow Z}{\Gamma \Rightarrow \bigcirc Z}
\qquad \text{barren},\quad \bigcirc Z \in \mathrm{Sf}^{R}(G)
\]
where \emph{barren} says the world this refutation builds declared no
modal successor. It is an index on the regular judgment, not a computed
property: a function on the family cannot be defined while the family is.
Obligation: \texttt{not\_force\_circ\_of\_no\_promise}, proved.

\paragraph{The promise join.} A join carries one additional regular
premise whose world becomes the modal successor of the new world, with
three side conditions: every modal formula kept in the conclusion has its
argument forced at the promise world; every modal formula kept is itself
forced there; and the whole conclusion context is forced there.
Obligation: \texttt{circ\_intro}, proved.

\subsection{The decision the screens force}
\label{sec:decision}

The two positive screens do not settle the design between them. A world
that forces everything witnesses every modal formula at once: enough for
the first screen, not for the second, which needs two worlds with
\emph{different} witnesses. Conversely the promise join alone cannot reach
every countermodel, and the sharp instance is $\bigcirc p \supset p$.

Its right subformulas are the goal and $p$, and contain no falsehood, so
the only regular axiom is the one at $p$ and both axioms at $p$ delete it:
no derivable regular context of this goal contains $p$, so no promise
premise can force $p$. Semantically the reason is sharper. The
countermodel needs a world forcing $p$; there the modality holds by its
unit, hence so does the goal, and falsehood is not a right subformula ---
so that world refutes \emph{nothing} available to the calculus, while
every world of the extracted model is a p-sequent and refutes its own goal
by construction.

So the witness cannot be a p-sequent. Three ways out, and only one
survives:

\begin{center}
\begin{tabular}{p{0.46\textwidth}p{0.44\textwidth}}
\hline
route & verdict\\
\hline
declare a context and assert a world forcing it &
\textbf{refuted} by the third screen: needs realisability, which is not
membership\\
let the extracted model carry worlds that are not p-sequents, labelled by
a context and refuting nothing &
the same obligation in different clothing --- such a label must be
realisable\\
a \emph{fallible} world &
no realisability obligation at all: a fallible world forces everything
\emph{by construction}\\
\hline
\end{tabular}
\end{center}

\noindent
The recommendation is therefore fallible worlds, with two consequences to
be accepted knowingly. The witness must cover the whole cone above the new
world, so the fallible world is a maximum of the model and no world of
such a model is barren --- correct rather than unfortunate, since a model
with a fallible top forces every modal formula everywhere. And it changes
the modality-free semantics: falsehood becomes forceable. The logic is
unchanged, and soundness is unaffected because every model the calculus
\emph{builds} can leave the fallible set empty unless a rule says
otherwise; but completeness starts from an \emph{arbitrary} countermodel,
which may now have fallible worlds, and the construction must handle them.
That cost falls on the next stage, not this one.

%% ==================================================================
\section{Status}
\label{sec:status}

\textsc{Proved}, sorry-free and pinned: everything in
\S\S\ref{sec:prelim}--\ref{sec:choice}, for modality-free goals; the
extension of \S\ref{sec:w1}; and the five semantic obligations and three
screens of \S\ref{sec:rules}.

\textsc{Open}: soundness and completeness for any calculus containing
modal rules. No modal rule is part of the calculus; the rules of
\S\ref{sec:rules} are put up for sign-off, and the decision of
\S\ref{sec:decision} is the substantive one, because it changes the model
structure and puts a new obligation on the completeness construction.

\appendix
%% ==================================================================
\section{The Lean development}
\label{app:code}

Every definition and proof discussed above appears below, unedited, in
dependency order. The library builds from scratch with no sorries, no
\texttt{native\_decide}, and the axiom pins of \S\ref{app:audit}.

\subsection{Preliminaries: language, models, forcing, closure}
\label{app:basic}
\lstinputlisting{../../FRJ/Basic.lean}

\subsection{The calculus}
\label{app:calculus}
\lstinputlisting{../../FRJ/Calculus.lean}

\subsection{The step relation and well-formedness}
\label{app:step}
\lstinputlisting{../../FRJ/Step.lean}

\subsection{The fresh-root construction}
\label{app:model}
\lstinputlisting{../../FRJ/Model.lean}

\subsection{Extraction of the model}
\label{app:extract}
\lstinputlisting{../../FRJ/Extract.lean}

\subsection{Soundness}
\label{app:soundness}
\lstinputlisting{../../FRJ/Sound.lean}

\subsection{The determining part}
\label{app:completelean}
\lstinputlisting{../../FRJ/Complete.lean}

\subsection{The completeness construction}
\label{app:minimal}
\lstinputlisting{../../FRJ/Minimal.lean}

\subsection{The modality: semantics, obligations, screens}
\label{app:modallean}
\lstinputlisting{../../FRJ/Modal.lean}

\subsection{The axiom audit}
\label{app:audit}
\lstinputlisting{../../FRJ/Audit.lean}

%% ==================================================================
\section{The archived parallel development}
\label{app:archive}

A parallel re-derivation was carried out before the modality-free
development was finished, and is superseded by it; it is archived rather
than deleted. Its modal semantics and screens have since been re-homed onto the live
development (Appendix~\ref{app:modallean}); what remains of interest here
is the parallel rule table, which reached index-equality by never
computing an index, and the replayed refutations from the article.

\subsection{Constraint models}
\lstinputlisting{../../Archive/FRJLax-parallel/Model.lean}

\subsection{The parallel rule table}
\lstinputlisting{../../Archive/FRJLax-parallel/Calculus.lean}

\subsection{The paper's own refutations, replayed}
\lstinputlisting{../../Archive/FRJLax-parallel/Paper.lean}

\subsection{The first modal refutation}
\lstinputlisting{../../Archive/FRJLax-parallel/Circ.lean}

\end{document}
